% Round 15. Supplied fixed-tree actual presentations; independently reviewed.
\section{An ambient whole-tree calculation on a fixed combinatorial chamber}
\label{r15ft:section}

A supplied affine family of one tree type admits a calculation which
retains the complete balanced cycle before intersecting the input.
This differs from substituting a positive-delta form into piecewise
affine open-ray maps. The hypotheses below are explicit presentation
data, not an assertion that every ordinary point admits them.

\begin{definition}[A split affine-tree presentation]
\label{r15ft:data}
Let $S/K$ be smooth integral of dimension $d$ and let
$U\subset S^{\rm an}$ be ordinary open. Let $b=(b_1,\ldots,b_N)$
be analytic units on $U$, and let $a_1,\ldots,a_s$ be regular
functions on a Zariski neighborhood, with distinct pairwise
differences which are units on $U$ and are included in $b$.
Put
\[
 W=(U\times\mathbf A^1)\setminus\bigcup_i\{z=a_i\},
       \qquad f=(b,z-a_1,\ldots,z-a_s).
\tag{FT15.1}\label{r15ft:moments}
\]
A split affine-tree presentation consists of an ambient open
$O\subset\mathbf R^N$ containing $t_b(U)$ and one finite embedded
labelled tree type with unbounded ends. Its vertices are affine
functions $q_v:O\to\mathbf R^s$. Its oriented primitive edge
directions are fixed; a bounded edge from $v$ to $w$ satisfies
\[
 q_w(u)-q_v(u)=\ell_{vw}(u)e_{vw},\qquad\ell_{vw}(u)>0,
\tag{FT15.2}\label{r15ft:edges}
\]
with affine $\ell_{vw}$. Every vertex is balanced. All relative
edge weights are one in the primitive disk-radius coordinate.
The resulting tree $\Gamma_u$ equals the full norm image of the
split marked fiber under $(v(z-a_i))_i$ for $u=t_b(x)$.
Bivalent auxiliary vertices are allowed.

The input has a supplied single presentation
\[
       \alpha=t_f^\star\omega,
       \qquad\omega\in B^{p,q}(O\times\mathbf R^s),
\tag{FT15.3}\label{r15ft:input}
\]
whose ambient support is proper over $O$. Over every compact
base set it is thus contained in a bounded vertical box. The
resulting ordinary zero extension over the omitted root and
infinity sections is part of this notation. The ambient form
is not assumed PWS or closed.
\end{definition}

The affine condition holds for supplied fixed split tree data on
a chamber with positive bounded-edge lengths, after choosing
affine extensions of its relations. It is a condition on an
ordinary neighborhood, not merely on one fiber. New edges opening
at the chosen point may prevent such a common chamber.

\begin{lemma}[Stars minus lines give a balanced ambient family]
\label{r15ft:family}
The data give an ambient $\delta$-cycle $\mathscr T$ on
$O\times\mathbf R^s$ of relative dimension one, with the
whole-family weighted support of the trees $\Gamma_u$.
\end{lemma}

\begin{proof}
For each vertex let $\operatorname{Star}_v$ be the full balanced
fan of its outgoing rays, translated to the origin. Transport
$\mathbf R^N\times\operatorname{Star}_v$ by the affine
automorphism $(u,w)\mapsto(u,q_v(u)+w)$ and restrict to $O$;
call this cycle $T_v$. For each bounded edge let $L_e$ be the
full affine line family in its direction through either endpoint.
All these objects are ambient $\delta$-cycles. Put
\[
                   \mathscr T=\sum_vT_v-\sum_{e\ {\rm bounded}}L_e.
\tag{FT15.4}\label{r15ft:star-line}
\]
Along a bounded edge the opposite rays from its endpoints cover
the intervening segment twice and each outside half-line once.
Subtracting the full line leaves precisely that segment. Every
unbounded ray occurs once. The same cancellation holds on the
joint base--fiber cells because the positions and lengths are
affine. Thus this is an equality of whole weighted cycles.
There is no restriction of a distribution to an isolated vertex
or a half-open edge, and later intersections retain vertex mass.
\end{proof}

\begin{theorem}[Actual projection for a supplied affine-tree presentation]
\label{r15ft:projection}
Under Definition~\ref{r15ft:data}, with
$\operatorname{pr}:O\times\mathbf R^s\to O$, set
\[
 \beta_0=\operatorname{pr}_*(\mathscr T\wedge\omega),
       \qquad\beta=t_b^\star\beta_0.
\tag{FT15.5}\label{r15ft:output}
\]
Then $\beta$ is an ordinary actual form of degree $(p-1,q-1)$
on $U$, with negative degrees interpreted as zero, and
\[
 \int_U\beta\wedge\eta
       =\int_{U\times\mathbf P^1}\alpha\wedge\pi^*\eta
\tag{FT15.6}\label{r15ft:comparison}
\]
for every ordinary smooth compact complementary test $\eta$.
All mixed tests and every positive-delta input stratum are
included. Different supplied presentations give the same actual
output on an ordinary overlap. All degrees are native over $K$.
\end{theorem}

\begin{proof}
The ambient affine $\delta$ calculus makes
$\mathscr T\wedge\omega$ a $\delta$-form with proper support
over $O$, and its affine pushforward is an ambient $\delta$-form
of the claimed total bidegree. Genuine unit pullback gives the
ordinary actual $\beta$. No closed ambient lifting is used.

For the prescribed pairing, choose an affinoid chart of the test
with support away from its analytic boundary, and append its
entire unit tuple $c$ to $b$. Write $b'=(b,c)$, and let
$l:\mathbf R^{N'}\to\mathbf R^N$ be coordinate projection.
Include also the unit moments of scalar localization cutoffs,
chosen to equal one near the paired support and zero near the
analytic chart boundary, as in Proposition~\ref{r15ray:analytic}.
On a common finite subdivision away from those boundary images,
the weighted source tropical cycle is
\[
 T_{\rm src}=\operatorname{pr}'{}^*T_{\rm base}\wedge
                         (l\times\mathrm{id})^*\mathscr T.
\tag{FT15.7}\label{r15ft:cycle-product}
\]
Here the first pullback is product with the full vertical space.
This equality is checked on maximal cells, before intersecting
with the positive-delta input. A point on a maximal source
skeletal cell has source moment graded transcendence degree $d+1$.
Let $B$ be the graded field generated over the graded ground field
by the initial forms of $b'$, and choose $i$ for which
$v(w_i)$ is maximal, where $w_i=z-a_i$.
Since $w_j=w_i+(a_i-a_j)$ and every nonzero difference occurs in $b$,
the ultrametric rule puts every initial $w_j$ in $B(\widetilde w_i)$.
At equal valuations, cancellation would give $v(w_j)>v(w_i)$,
contradicting the choice of $i$. Thus the full source moment field
has transcendence degree at most $\operatorname{trdeg}B+1$.
It follows that $\operatorname{trdeg}B=d$; this proves full rank
for the base moment tuple, not merely for the completed residue
field of the base point.
The latter graded field is algebraic over $B$, so
$\widetilde w_i$ remains transcendental over it. The polynomial
Gauss-norm criterion now identifies the relative point as a
Gauss point of the split marked tree.
Conversely these Gauss disks over full-rank base moment points
realize every base maximal cell and every relative edge.
A finite coordinate map on the base, combined with the primitive
radius coordinate $z-a_i$, has the same degree as that base map.
The native source weight is therefore the base weight times the
relative edge weight one. This is the calibrated degree argument
of Proposition~\ref{r15ray:analytic}, used to identify a whole tropical cycle.
The affine family in \eqref{r15ft:star-line} has exactly those
maximal cells and weights, proving \eqref{r15ft:cycle-product}.
A vertex-supported input has not been discarded: it has not yet
been intersected with this cycle.

If $\eta_0$ is the smooth ambient test presentation, integration
of actual top forms and the ambient affine projection formula give
\[
\begin{split}
 \int T_{\rm src}\wedge(l\times\mathrm{id})^*\omega
                                \wedge\operatorname{pr}'{}^*\eta_0
 &=\int T_{\rm base}\wedge
       \operatorname{pr}'_*(l\times\mathrm{id})^*
                       (\mathscr T\wedge\omega)\wedge\eta_0\\
 &=\int T_{\rm base}\wedge l^*\beta_0\wedge\eta_0.
\end{split}
\tag{FT15.8}\label{r15ft:refined-calculation}
\]
The base change here is only Cartesian product by the extra
Euclidean test coordinates: $l$ is a surjective coordinate
projection. It follows from affine current pushforward and flat
product pullback, or from cellwise Fubini. No pullback of arbitrary
analytic currents is asserted. The last integral is the pairing
with $\beta$. A partition on the compact test support proves
\eqref{r15ft:comparison}. Actual-current injectivity gives
uniqueness and ordinary gluing.
\end{proof}

\begin{corollary}[Two finite marks and native finite descent]
\label{r15ft:two-marks}
With two distinct finite marks and infinity, the tree is a
translated three-ray fan with vertex
$(v(a_1-a_2),v(a_1-a_2))$. This is affine in an included base
moment. Hence every supplied proper ambient actual presentation
using those fiber units has the stated ordinary output, without
any bounded-edge chamber condition. One finite mark gives the
straight-line version.

If such presentations, or general fixed-tree ones, are available
on an ordinary covering of the full inverse image under a finite
normal algebraic splitting cover to which the smooth-base presentation
hypotheses apply (in particular a finite etale cover over the working
open), the outputs descend whenever
the output degree is in the established actual finite-trace strip.
In particular input $(2,2)$ has output $(1,1)$ in that strip.
\end{corollary}

\begin{proof}
The three primitive directions are $(1,0),(0,1),(-1,-1)$, with
weights one, so the affine-family construction applies. For finite
descent, glue upstairs by the prescribed current and apply the
frozen actual trace in the output strip, divided by the full
native degree. On every downstairs smooth test, the source finite
product cover contributes exactly that same degree. This gives
the original current. No trace in a higher actual degree is used.
\end{proof}

\paragraph{Why this fixed-tree calculation alone does not cross every wall.}
The theorem needs a supplied ambient input and a common affine
whole-tree family. It does not apply a delta pullback to a
piecewise affine ray map. The subsequent universal graphic-matroid
kernel in Section~\ref{r15ug:section} supplies the whole-cycle
identity on every test refinement across changing tree types.
The ambient atlas in Section~\ref{r15ad:section} and the available
finite output trace then give Theorem~\ref{r15a2:projection}.
The present fixed-tree construction remains a direct local proof,
including the moving-vertex example; balancedness of a pushed
current by itself is never substituted for an ambient representative.
